% Part: computability
% Chapter: computability-theory
% Section: s-m-n

\documentclass[../../../include/open-logic-section]{subfiles}

\begin{document}

\olfileid{cmp}{thy}{smn}
\olsection{Teorema $s$-$m$-$n$}

\begin{explain}
Teorema sabanjure diarani ``teorema $s$-$m$-$n$'',
kanthi alesan kang sedhela maneh bakal cetha.
Perkara kang angel yaiku mangerteni persis
apa kang diandharake teorema iki; yen
andharane wis dingerteni, teorema iki
bakal katon cukup cetha.
\end{explain}

\begin{thm}
\ollabel{thm:s-m-n}
  Kanggo saben pasangan wilangan asli $n$ lan~$m$,
  ana fungsi rekursif primitif~$s^m_n$ kang
  ndadekake, kanggo saben rerangken
  $e$, $a_0$, \dots, $a_{m-1}$, $y_0$, \dots,
  $y_{n-1}$, kita nduwé
  \[
  \cfind{s^m_n(e, a_0, \dots, a_{m-1})}[n](y_0, \dots, y_{n-1}) \simeq
  \cfind{e}[m+n](a_0, \dots, a_{m-1}, y_0, \dots, y_{n-1}).
\]
Cathetan panyunting: pangindeksan ing teorema iki yaiku pangindeksan program baku kang ngidini konstruksi kode rekursif primitif, ora sembarang enumerasi universal. Yen aritas nol diidinake, tuple input utawa parameter kang cocog kosong; yen ora, pratelan diwaca kanggo aritas kang positif.
\end{thm}

\begin{explain}
Migunani yen $s^m_n$ dianggep tumindak marang
\emph{program}. Tegese, $s^m_n$ nampa program~$e$
kanggo fungsi aritas $(m+n)$, uga input tetep
$a_0$, \dots, $a_{m-1}$; banjur ngasilake program
% OLPL-133: Sumber beku nulis x ing rong rujukan indeks program ing paragraf iki; teorema lan ukara sadurunge nggunakake e.
$s^m_n(e, a_0, \dots, a_{m-1})$ kanggo fungsi
aritas $n$ saka argumen kang isih kari. Cathetan koreksi sumber OLPL-133: rong rujukan program x kang ora ditegesi ing paragraf Inggris dibenerake dadi e, manut teorema lan input program kang disebut ing ukara sadurunge.
\iftag{TMs}{Yen $e$ dianggep minangka jlentrehan
mesin Turing, mula $s^m_n(e, a_0, \dots, a_{m-1})$
iku mesin Turing kang, yen nampa input $y_0$,
\dots,~$y_{n-1}$, nambahake $a_0$, \dots,~$a_{m-1}$
ing wiwitan rerangken input, banjur nglakokake~$e$.
Saben $s^m_n$ dadi fungsi rekursif primitif
kanggo nemokake kode mesin Turing kang cocog.}{}
\end{explain}

\end{document}
